\documentclass[spanish,letterpaper,oneside]{article}
\def\documenttitle {Antecedentes historicos de los derechos humanos}
\def\documentsubtitle {Cuestionario diagnostico y fundamentos historicos}
\def\documentsubject {Actividad 1 - Antecedentes de los derechos humanos}
\def\documentauthor {Martin Jonathan de la Cruz}
\def\coursename {Antecedentes de los derechos humanos}
\def\coursecode {005}
\def\universityname {Universidad Abierta y a Distancia de Mexico}
\def\universityfaculty {Licenciatura en Derecho}
\def\universitydepartment {Antecedentes de los derechos humanos}
\def\universitydepartmentimage {departamentos/UnADM}
\def\universitydepartmentimagecfg {height=1.57cm}
\def\universitylocation {Roma Norte, Ciudad de Mexico}
\def\authortable {\begin{tabular}{ll}
Alumno: & Martin Jonathan de la Cruz \\
Matricula: & ES2611202040 \\
Actividad interna: & 1 \\
Semana: & 1 \\
Figura docente: & Daniel Alexis Lozano Ortega \\
Fecha: & \today
\end{tabular}}
\input{template}
\usepackage{helvet}
\renewcommand{\familydefault}{\sfdefault}
\setcitestyle{authoryear,open={(},close={)}}
\begin{document}
\templatePortrait
\templatePagecfg
\templateFinalcfg
\onehalfspacing
\section{Introduccion}
El diagnostico inicial permite reconocer conocimientos previos para
orientar el estudio y familiarizarse con las evaluaciones del Aula
Virtual. Los tres reactivos relacionan limitacion del poder monarquico,
declaraciones revolucionarias e internacionalizacion de los derechos.
No representan una evaluacion definitiva del conocimiento juridico.
Las respuestas proceden de la revision del intento proporcionada por
el estudiante; las justificaciones siguientes se desarrollaron despues
mediante contraste documental, no son explicaciones escritas durante
el examen.

\section{Cuestionario sobre antecedentes y declaraciones}
\subsection{Limitacion del poder monarquico}
\textbf{Pregunta 1.} ¿Cual de los siguientes documentos es considerado
un antecedente historico relevante de la limitacion del poder del
monarca en Inglaterra?
\begin{enumerate}
\item[a)] La Declaracion Universal de los Derechos Humanos.
\item[b)] La Carta Magna de 1215.
\item[c)] La Declaracion de Independencia de los Estados Unidos.
\item[d)] La Constitucion mexicana de 1917.
\end{enumerate}
\textbf{Respuesta:} b), la Carta Magna de 1215. La revision la registra
como correcta.

\textbf{Justificacion:} la Carta Magna establecio limites al poder
real en el contexto de una crisis entre el rey y los barones ingleses.
Su relevancia consiste en someter determinadas actuaciones del monarca
a reglas; no debe confundirse con una declaracion universal de derechos,
pues buscaba proteger intereses y libertades de grupos concretos
\citep{parlamentoCarta1215ADH}.

\subsection{Declaracion francesa de 1789}
\textbf{Pregunta 2.} ¿Que acontecimiento historico se relaciona
directamente con la Declaracion de los Derechos del Hombre y del
Ciudadano de 1789?
\begin{enumerate}
\item[a)] La Revolucion Industrial.
\item[b)] La Primera Guerra Mundial.
\item[c)] La Revolucion Francesa.
\item[d)] La Independencia de Mexico.
\end{enumerate}
\textbf{Respuesta:} c), la Revolucion Francesa. La revision la registra
como correcta.

\textbf{Justificacion:} la Asamblea Nacional francesa adopto la
Declaracion el 26 de agosto de 1789 en el proceso revolucionario.
Sus articulos 1 y 3 reconocen igualdad juridica y soberania nacional;
el articulo 16 relaciona constitucion, garantia de derechos y separacion
de poderes. Estos contenidos explican su vinculacion con el cambio
politico frances, no con los otros acontecimientos propuestos
\citep[arts.~1, 3 y 16]{declaracionFrancia1789ADH}.

\subsection{Internacionalizacion de los derechos humanos}
\textbf{Pregunta 3.} La Declaracion Universal de los Derechos Humanos
fue adoptada por la Organizacion de las Naciones Unidas en 1948,
despues de la Segunda Guerra Mundial.

\textbf{Opciones:} Verdadero o Falso.

\textbf{Respuesta:} Verdadero. El texto compartido identifica esta
como la respuesta apropiada y contiene retroalimentacion afirmativa;
no se dispone de la marca visual original de seleccion.

\textbf{Justificacion:} la Asamblea General de las Naciones Unidas
proclamo la Declaracion el 10 de diciembre de 1948 en Paris, mediante
la resolucion 217 A (III), en el contexto posterior a la guerra.
Establecio un ideal comun de proteccion para todos los pueblos y
naciones. Es una declaracion, no un tratado, y su proclamacion no
demuestra por si sola cumplimiento efectivo \citep{onuDeclaracion1948ADH}.

Los tres referentes muestran cambios en el alcance de la proteccion:
limites pactados al monarca, reconocimiento revolucionario de derechos
y un parametro internacional. Esa trayectoria admite avances y
retrocesos, no una expansion automatica ni uniforme de garantias
\citep[p.~11]{alvarezIcaza2009ADH}\citep{unedVideo2024}.

\clearpage
\section{Conclusiones}
Considero importante distinguir un antecedente historico de una garantia
universal: limitar al rey en 1215 no equivale a reconocer los mismos
derechos a toda persona. La Declaracion francesa y la Declaracion
Universal amplian el marco de comprension, pero su existencia exige
examinar tambien las condiciones reales de proteccion.

La revision registra el intento terminado y los reactivos no calificados.
No corresponde convertir las respuestas correctas en una nota de
10/10 ni afirmar dominio completo de la asignatura. Para mi formacion,
el contraste posterior sirve para justificar las respuestas y distinguir
fechas, contextos y alcances, sin atribuir al momento del examen las
fuentes consultadas despues.\footnote{GitHub Copilot apoyo la redaccion
y el contraste documental posterior. No respondio ni envio este intento
del cuestionario.}

\clearpage
\begin{thebibliography}{5}
\bibitem[Alvarez Icaza Longoria(2009)]{alvarezIcaza2009ADH}
Alvarez Icaza Longoria, E. (2009). \emph{Los derechos humanos en Mexico}.
Nostra Ediciones.
\bibitem[Asamblea Nacional Constituyente de Francia(1789)]{declaracionFrancia1789ADH}
Asamblea Nacional Constituyente de Francia. (1789, 26 de agosto).
\emph{Declaration of Human and Civic Rights of 26 August 1789}.
Conseil constitutionnel.
\url{https://www.conseil-constitutionnel.fr/en/declaration-of-human-and-civic-rights-of-26-august-1789}
\bibitem[Naciones Unidas(1948)]{onuDeclaracion1948ADH}
Naciones Unidas. (1948, 10 de diciembre).
\emph{Declaracion Universal de Derechos Humanos}.
\url{https://www.un.org/en/about-us/universal-declaration-of-human-rights}
\bibitem[UK Parliament(s.\,f.)]{parlamentoCarta1215ADH}
UK Parliament. (s.\,f.). \emph{Magna Carta}.
\url{https://www.parliament.uk/about/living-heritage/evolutionofparliament/originsofparliament/birthofparliament/overview/magnacarta/}
\bibitem[Universidad Nacional de Educacion a Distancia(2024)]{unedVideo2024}
Universidad Nacional de Educacion a Distancia. (2024, 23 de octubre).
\emph{Tema 11. Los Derechos Humanos Concepto y Fundamentacion} [Video].
Canal UNED. \url{https://canal.uned.es/video/671b6728e0637d00c70d7b80}
\end{thebibliography}
\end{document}
